\documentclass[../../../include/open-logic-section]{subfiles}

\begin{document}

\olfileid{sth}{replacement}{absinf}
\olsection{جایگزینی و «نامتناهیِ مطلق»}

اکنون~\limofsize{} را کنار می‌گذاریم و خانوادهٔ دیگری از کوشش‌های
(درونی) برای توجیه جایگزینی را بررسی می‌کنیم که اندیشهٔ تشکیل مجموعه‌ها
در مراحل را جدی می‌گیرند.

هنگامی که نخستین بار فرایند تکراری را ترسیم کردیم، اصولی برای توضیح آنچه
در هر مرحله رخ می‌دهد عرضه کردیم. آن‌ها~\stageshier، \stagesord و
\stagesacc بودند. بعدتر اصولی افزودیم که دربارهٔ شمار مراحل اطلاعاتی به
ما می‌دادند: \stagessucc{} می‌گفت فرایند تشکیل مجموعه هرگز پایان
نمی‌یابد، و~\stagesinf{} می‌گفت فرایند از مرحله‌ای نامتناهی می‌گذرد.

معقول است پیشنهاد کنیم این دو اصل اخیر از اصلی گسترده‌تر مانند اصل زیر
نتیجه می‌شوند:
\begin{enumerate}
	\item[] \stagesinex. مطلقاً بی‌نهایت مرحله وجود دارد؛ سلسله‌مراتب تا
	حد ممکن بلند است.
\end{enumerate}
بدیهی است که این اصلی غیرصوری است. اما حتی اگر بی‌درنگ از تصور
انباشتی-تکراریِ مجموعه \emph{نتیجه نشود}، بی‌گمان با آن \emph{همساز} به
نظر می‌رسد. دست‌کم، برخلاف~\limofsize، این اندیشه را حفظ می‌کند که
مجموعه‌ها مرحله‌به‌مرحله تشکیل می‌شوند.

اکنون امید آن است که از~\stagesinex{} برای توجیه جایگزینی بهره بگیریم.
پس ببینیم این کار چگونه ممکن است.

در~\olref[ordinals][idea]{sec} دیدیم ساختن خوش‌ترتیبی‌ای که (از نظر
مقصود ریاضی) باید با~$\omega+\omega$ هم‌ریخت باشد آسان است. به‌بیانی
دیگر، به‌آسانی می‌توانیم فرایندی تکراری و مرحله‌به‌مرحله را تصور کنیم که
نوع ترتیبش (از نظر مقصود ریاضی)~$\omega+\omega$ باشد. پس اگر
\stagesinex{} را پذیرفته باشیم، بی‌گمان باید بپذیریم دست‌کم مرحلهٔ
$\omega+\omega$ام سلسله‌مراتب، یعنی~$V_{\omega+\omega}$، وجود دارد؛
زیرا سلسله‌مراتب قطعاً \emph{می‌توانست} تا اینجا ادامه یابد.

این اندیشه چنین تعمیم می‌یابد: برای هر خوش‌ترتیب، فرایند ساختن
سلسله‌مراتب تکراری باید دست‌کم به‌اندازهٔ آن خوش‌ترتیب پیش رود. و می‌توانیم
این امر را تنها با اصل موضوع انگاشتنِ
\olref[ordinals][ordtype]{thmOrdinalRepresentation} تضمین کنیم. این اصل
به ما می‌گوید هر خوش‌ترتیب با یک عدد ترتیبی فون نویمان هم‌ریخت است. چون
هر عدد ترتیبی فون نویمان با رتبهٔ خودش برابر خواهد بود،
\olref[ordinals][ordtype]{thmOrdinalRepresentation} به ما خواهد گفت
هرگاه بتوانیم در نظریهٔ مجموعه‌های خود خوش‌ترتیبی‌ای را توصیف کنیم،
فرایند تکراریِ ساختن مجموعه باید از آن خوش‌ترتیب فراتر رود.

این اندیشه بی‌گمان پیامدی از~\stagesinex{} به نظر می‌رسد. متأسفانه اگر
هدفمان استخراج جایگزینی از این اندیشه باشد، با مانعی ساده و فنی
روبه‌رو می‌شویم: جایگزینی اکیداً قوی‌تر از
\olref[ordinals][ordtype]{thmOrdinalRepresentation} است. (این نکته را
\citet[\S13.2]{Potter2004} مطرح کرده است؛ ما آن را در
\olref[sth][replacement][finiteaxiomatizability]{sec} اثبات خواهیم کرد.)

نتیجه آن است که اگر بخواهیم~\stagesinex{} را چنان بفهمیم که جایگزینی
را به دست دهد، نمی‌تواند \emph{صرفاً} بگوید سلسله‌مراتب از هر خوش‌ترتیبی
فراتر می‌رود؛ باید ادعایی قوی‌تر مطرح کند. برای این منظور،
\citet{Shoenfield:AST} تقویت بسیار طبیعیِ زیر را پیشنهاد کرد: سلسله‌مراتب
با هیچ مجموعه‌ای \emph{هم‌نهایی} نیست.\footnote{به نظر می‌رسد گودل
اندیشه‌ای مشابه پیشنهاد کرده باشد؛ بنگرید به
\citet[p.~223]{Potter2004}. برای بحثی دربارهٔ گودل و
\citeauthor{Shoenfield:AST}، بنگرید به~\citet[90--5]{Incurvati2020}.}
به‌تفصیل بیشتر: اگر~$\tau$ نگاشتی باشد که مجموعه‌ها را به مراحل
سلسله‌مراتب می‌فرستد، تصویر هیچ مجموعهٔ~$A$ تحت~$\tau$ سلسله‌مراتب را
تمام نمی‌کند. به‌بیانی دیگر (به‌صورت طرحواره‌ای):
\begin{enumerate}
	\item[] \stagescofin. اگر~$A$ یک مجموعه باشد و برای هر~$x\in A$،
	$\tau(x)$ یک مرحله باشد، آنگاه مرحله‌ای وجود دارد که پس از هر
	$\tau(x)$ برای~$x\in A$ می‌آید.
\end{enumerate}
بدیهی است که~$\ZF$ صورت‌بندی مناسبی از~\stagescofin{} را اثبات می‌کند.
در جهت عکس، می‌توانیم به‌طور غیرصوری استدلال کنیم که~\stagescofin{}
جایگزینی را توجیه می‌کند.\footnote{اثبات جایگزینی با استفاده از
صورت‌بندی صوریِ~\stagescofin{} دشوارتر خواهد بود، زیرا~$\Z$ به‌تنهایی
برای تعریف مراحل به‌اندازهٔ کافی قوی نیست و ازاین‌رو روشن نیست
\stagescofin{} را چگونه باید صوری کرد. یک گزینه آن است که در امتدادی از
$\LT$ کار کنیم، چنان‌که در~\olref[sth][replacement][extrinsic]{sec}
بحث شد.} زیرا فرض کنید
$(\forall x \in A)\lexists![y][\phi(x,y)]$. برای هر~$x\in A$، بگذارید
$\sigma(x)$ همان~$y$ای باشد که~$\phi(x,y)$، و بگذارید~$\tau(x)$ مرحله‌ای
باشد که~$\sigma(x)$ نخستین بار در آن تشکیل می‌شود. بنا بر~\stagescofin،
مرحله‌ای مانند~$V$ وجود دارد به‌طوری‌که
$(\forall x\in A)\tau(x)\in V$. چون~$V$ یک مرحله و بنابراین توانمند است،
از~$\tau(x)\in V$ و~$\sigma(x)\subseteq\tau(x)$ نتیجه می‌شود
$\sigma(x)\in V$. اکنون با جداسازی می‌توانیم مجموعهٔ
$\Setabs{y \in V}{(\exists x \in A)\sigma(x) = y} =
\Setabs{y}{(\exists x \in A)\phi(x,y)}$ را به دست آوریم.

\begin{prob}
	\stagescofin{} را درون~$\ZF$ صوری کنید.
\end{prob}

پس~\stagescofin{} جایگزینی را موجه می‌سازد. و دست‌کم پذیرفتنی است که
\stagesinex{} نیز~\stagescofin{} را موجه سازد. زیرا فرض کنید
\stagescofin{} شکست بخورد. در این صورت سلسله‌مراتب با مجموعه‌ای مانند~$A$
هم‌نهایی است؛ یعنی نگاشتی مانند~$\tau$ داریم به‌طوری‌که برای هر مرحلهٔ~$S$
یک~$x\in A$ وجود دارد که~$S\in\tau(x)$. در آن صورت راهی داریم تا
«نامتناهیِ مطلق» مفروضِ سلسله‌مراتب را مهار کنیم: این نامتناهی با برد
$\tau$ روی~$A$ \emph{تمام می‌شود}. و این امر اندیشهٔ «مطلقاً نامتناهی»
بودن سلسله‌مراتب را تضعیف می‌کند. با گرفتن عکس نقیض:
\stagesinex{} مستلزم~\stagescofin{} است، و~\stagescofin{} نیز به‌نوبهٔ
خود جایگزینی را توجیه می‌کند.

این استدلال کوششی واقعاً امیدبخش برای ارائهٔ توجیهی درونی از جایگزینی
است. اما اینکه در نهایت موفق می‌شود یا نه، پرسشی است که باید داوری‌اش را
به شما واگذار کنیم.

\end{document}
